\begin{enumerate}[label=\textbf{\ansnum.\arabic*}]
  \item $0.005\cdot4=0.02$\,s, so $50$\,Hz. For $50$\,Hz with the timestep $0.004$\,s: decimation $5$ ($5\cdot0.004=0.02$\,s).
  \item $0.3\cdot10=3$\,N. An action of $-4$ is clamped to $-1$, giving $-10$\,N. The clamp is safe because an exploring policy with a wide Gaussian noise often outputs values far outside the range; without the clamp these would be huge forces.
  \item $0.9\cdot0.8\cdot0.95\cdot0.99=0.677$. With $0.1$ for the first factor: $0.1\cdot0.8\cdot0.95\cdot0.99=0.075$. The average of the four factors is $0.91$ in the first case and $(0.1+0.8+0.95+0.99)/4=0.71$ in the second. The product drops to $11\%$ of its value when one factor is bad; the average only drops by a fifth. A product reward punishes a failure in \emph{any} factor strongly.
  \item $\theta=\pi+0.1$: $\cos=-\cos0.1=-0.995$, $\sin=-\sin0.1=-0.0998$. $\theta=-\pi+0.1$: $\cos=-0.995$, $\sin(-\pi+0.1)=-\sin(\pi-0.1)=-0.0998$, the same. $\theta=3\pi+0.1$ again the same pair. The raw angles $3.24$, $-3.04$ and $9.52$ are very different numbers for the same physical state, and the network would have to learn that they mean the same thing.
  \item $10\cdot20=200$ steps. $4096\cdot24=98\,304$ samples per iteration. The number of steps of an environment is $24k$ after $k$ iterations; a whole number of episodes of $200$ steps needs $24k$ to be a multiple of $200$: the smallest $k$ is $25$ ($600$ steps, $3$ episodes).
  \item At a \emph{failure} the episode ends because the future reward is zero (or very bad), so the target of the value function is the reward alone. At a \emph{time-out} the robot would have continued to collect reward, so the value of the last state is \emph{bootstrapped}: the target contains $\gamma V(s')$. If a time-out were treated as a failure, the value function would learn that the states shortly before the end of the episode are bad, for no reason in the state itself, and the policy would learn a wrong, time-dependent behaviour.
  \item The critic is not used on the real robot: it only judges the states during training, so it may use privileged information that makes its estimates more precise. Examples: the true linear velocity of the body (the K1 has no sensor for it, only an estimate), the contact forces of the feet, the height of the terrain, the friction coefficient and the mass that were randomised.
  \item I ran it. With the reward \emph{upright only} the training gives $K=(0.51, 3.88, 5.94, 6.04)$ and a mean return of $199.9$ (of $200$); the cart still ends close to the middle (mean $|x|=0.075$\,m, against $0.0076$\,m for the full reward) and no start terminates. The cart is not punished for being off-centre, so it is less precise, but within $10$\,s it does not run away. A longer episode would show more drift. (One seed.)
  \item I ran it: initial angle uniform in $\pm0.5$\,rad (the rest as in the chapter, $61$ iterations, seed $0$). The final policy $K=(0.84,2.61,6.00,4.21)$ has a mean return of $190.2$ (minimum $167.6$ over $1000$ starts), and $0\%$ of $2000$ starts are terminated. The linear policy still works; the return is lower because the pole starts farther from the upright and needs time to get back.
  \item I ran it (one seed, $2000$ evaluation starts, $61$ iterations): trained without noise and pushes, the final return is $199.3$ in a clean test and $177.1$ in a test with noise and pushes. Trained \emph{with} noise and pushes, it is $166.0$ in the clean test and $173.2$ in the noisy test. So for this tiny linear policy the training with noise did not help (the noisy-test return is $173.2$ against $177.1$, which is within what one seed can give), and the gains became a little larger, which made the clean test worse. The lesson: randomisation is not a free gain; a policy that has too few parameters to profit from it pays for it. It is the structure of the experiment that carries over to mjlab, not the numbers.
\end{enumerate}
